\chapter{Related Work}\label{ch:related}

This chapter surveys the work AGR builds on and competes with. It proceeds from
static graph-augmented generation, through path-retrieval and agentic
exploration, to verification and factuality measurement, and closes with a
structured comparison of the six closest systems. That comparison serves two
purposes: it locates the gap this thesis addresses, and it justifies the choice
of baselines in \Cref{sec:baselines}.

\section{Static Graph-Augmented Generation}\label{sec:static-graphrag}

\subsection{GraphRAG and Query-Focused Summarization}

GraphRAG~\cite{edge2024local} constructs a knowledge graph from a text corpus
using a language model, partitions it into communities, and pre-computes
summaries at each level, so a query can be answered from structured context
rather than retrieved passages. Its strength is global questions --- those whose
answer depends on the corpus as a whole rather than any single passage --- which
flat retrieval handles poorly.

The relevant limitation is that retrieval remains a single step preceding
generation. The context is assembled before the model has begun reasoning, so it
cannot be conditioned on intermediate findings. GraphRAG also builds its graph by
language-model extraction, which is appropriate for unstructured corpora but
introduces extraction error into the reference itself --- a consideration that
shapes the environment design in \Cref{sec:env-goals}.

\subsection{HippoRAG and Memory-Structured Retrieval}

HippoRAG~\cite{gutierrez2024hipporag} draws on the hippocampal indexing theory of
human memory, combining a knowledge graph with a personalised PageRank retrieval
step so that a single query can integrate evidence distributed across many
passages. It substantially improves multi-hop retrieval over flat baselines at
low cost.

Its retrieval is nonetheless a single graph-propagation step, not an
agent-directed traversal: the propagation is fixed and the model does not choose
where to look next. Temporal extensions such as T-GRAG~\cite{li2025tgrag} address
a different axis --- conflicting and redundant facts across time --- which is
orthogonal to the navigation question and is noted here for completeness rather
than compared against.

\section{Path-Retrieval and Semantic-Parsing KGQA}\label{sec:path-kgqa}

\subsection{Reasoning-on-Graphs}

Reasoning-on-Graphs (RoG)~\cite{luo2024reasoning} adopts a
planning--retrieval--reasoning pipeline: a fine-tuned model generates relation
paths grounded in the graph schema, those paths are used to retrieve matching
instantiations, and the model reasons over the retrieved paths. Because answers
are produced from retrieved paths, RoG is faithful by construction and can
present the path as an explanation.

Two properties distinguish it from the present work. It is \emph{fine-tuned},
which places its reported numbers outside direct comparison with prompting
methods. And its plan is generated in one shot before retrieval: if the generated
relation path does not instantiate in the graph, there is no mechanism to revise
it mid-course. The pre-generation verification question does not arise, because
grounding is guaranteed structurally rather than checked.

\subsection{StructGPT and KG-Agent}

StructGPT~\cite{jiang2023structgpt} provides a general interface for reasoning
over structured data --- knowledge graphs, tables, and databases --- through an
invoke--linearise--generate loop, in which specialised interfaces extract
evidence that is linearised into text for the model. It is deliberately general
rather than graph-specific, and performs no open-ended graph search.

KG-Agent~\cite{jiang2024kgagent} is closer to the agentic setting: a
multifunctional toolbox, a knowledge-graph executor, and a dynamic memory, driven
by an instruction-tuned model that autonomously selects tools across iterations.
It demonstrates that a small fine-tuned model with good tool abstractions is
competitive with much larger prompted ones. It plans implicitly through tool
selection rather than committing to explicit sub-objectives, provides no
backtracking mechanism, and performs no verification of the final answer.

\section{Agentic Graph Exploration}\label{sec:agentic-exploration}

\subsection{Think-on-Graph}

Think-on-Graph (ToG)~\cite{sun2024think} established the template for
training-free agentic KGQA. The model performs iterative beam search over the
graph: at each step it is shown candidate relations and entities and prunes them
itself, with beam width and maximum depth both set to three. Exploration
terminates when the model judges the retrieved paths sufficient, at which point
it answers directly.

ToG is the most directly comparable prior system to AGR --- training-free,
prompting-based, operating on Freebase, evaluated on the same benchmarks --- and
is therefore the agentic baseline in \Cref{sec:baselines}. Its limitations are
instructive. There is no explicit decomposition, so multi-constraint questions
are handled implicitly. Low-scoring branches are pruned within the beam, but
there is no mechanism to return to an abandoned node once the beam has moved on.
And the transition from exploration to answering is a single model judgement:
nothing checks the answer that follows.

\subsection{Plan-on-Graph}

Plan-on-Graph (PoG)~\cite{chen2024plan} is the closest system to AGR in ambition.
It decomposes the question into sub-objectives including constraint conditions,
explores adaptively with a breadth that varies by sub-objective, and adds
Guidance, Memory, and Reflection mechanisms. Reflection decides both whether to
self-correct and which entity to backtrack to, making PoG the one surveyed system
with genuine backtracking.

PoG's self-correction operates on the \emph{search process}: it reconsiders where
to explore. It does not decompose the draft answer into claims and check each
against the retrieved triples before responding. The distinction is the gap this
thesis occupies, and it is stated precisely in \Cref{sec:positioning}.

A naming hazard must be recorded here. A second, unrelated system is also
abbreviated PoG --- Paths-over-Graph~\cite{tan2024paths} --- and reports
considerably higher figures. Survey tables that cite ``PoG'' without
disambiguation are a known source of contaminated comparisons, and this thesis
always writes Plan-on-Graph or Paths-over-Graph in full.

\subsection{Generate-on-Graph and Reasoning with Trees}

Generate-on-Graph (GoG)~\cite{xu2024generate} targets the incomplete-graph
setting through a thinking--searching--generating loop: when the graph lacks a
required fact, the model generates candidate triples from its parametric
knowledge and continues. It also expands subgraphs dynamically to handle the
mediator nodes described in \Cref{sec:kgs}, which ToG tends to skip. The
trade-off is directly opposed to this thesis's: GoG deliberately admits
model-generated facts into the evidence, whereas AGR admits only graph triples.

Reasoning with Trees~\cite{shen2025reasoning} applies explicit tree search to
knowledge graph question answering. It is cited in \Cref{sec:mcts-clarification}
as an example of genuine tree-search machinery, in contrast with the guided
best-first search implemented here.

\subsection{KBQA-o1 and Genuine Tree Search}\label{sec:kbqao1}

KBQA-o1~\cite{luo2025kbqao1} is the closest system to the search strategy the
present work's proposal originally described, and therefore the most important
one to distinguish from it. It performs agentic knowledge base question
answering with full Monte Carlo Tree Search --- a policy model proposing actions,
a reward model scoring them, and rollouts whose values propagate back through the
tree --- wrapped around a ReAct-style agent that generates logical forms
step by step against a Freebase environment. The exploration additionally
generates annotations used for incremental fine-tuning, so the search improves
the model that drives it.

Two aspects bear on this thesis. First, KBQA-o1 is what the withdrawn
``MCTS-inspired'' description of \Cref{sec:mcts-clarification} would have
claimed: it has the rollouts and value backpropagation that AGR does not, and
naming the difference honestly is what keeps the two separable.

Second, its ablation is directly relevant to the methodological argument of
\Cref{sec:positioning}. Removing MCTS from KBQA-o1 while leaving every other
component intact drops overall GrailQA F1 from 78.5 to 48.5 --- a
thirty-point attribution to a single mechanism, obtained precisely because the
authors ablated components rather than reporting the assembly as a whole. That
is the same instrument this thesis applies in \Cref{sec:ablations}, and it is
evidence that the instrument is informative.

KBQA-o1 is not a baseline here. It is fine-tuned rather than prompted, evaluates
under a low-resource few-shot protocol, and does not report
ComplexWebQuestions, so no controlled comparison against it is available in this
environment.

\section{Verification and Self-Correction in Language
  Models}\label{sec:verification-related}

\subsection{Chain-of-Verification}

Chain-of-Verification (CoVe)~\cite{dhuliawala2024chain} has the model draft a
response, plan verification questions about it, answer those independently, and
regenerate a corrected response. It establishes the core premise of this thesis's
verification layer --- that decomposing a draft and checking its parts reduces
factual error --- but performs the checks against the model's own knowledge
rather than an external symbolic source. Multi-agent debate~\cite{du2024improving}
achieves a related effect by having several model instances critique one another.

This distinction matters given evidence that language models cannot reliably
self-correct reasoning without external feedback~\cite{huang2024selfcorrect}. A
verification step whose ground truth is the same model that produced the error is
limited by construction. AGR's verification is grounded in traversed triples,
making the check external to the generator.

\subsection{Ontology-Guided and Self-Correcting Graph RAG}

A self-correcting agentic graph RAG system for clinical decision support in
hepatology~\cite{hu2025selfcorrecting} demonstrates the combination of agentic
retrieval with a correction loop in a high-stakes domain, and ontology-guided
reasoning constrains an agent to logically valid relation use. These
substantiate the practical motivation for verification, though in domain-specific
settings rather than on the open-domain benchmarks used here. Reflexion-style
verbal self-improvement~\cite{shinn2023reflexion} provides the general agentic
pattern of which these are graph-specific instances.

\section{Measuring Factuality}\label{sec:factuality-related}

\subsection{Claim-Decomposition Metrics}

FActScore~\cite{min2023factscore} evaluates long-form generation by decomposing
text into atomic facts and verifying each against a reference source, reporting
the supported proportion. This decomposition-then-verify protocol is the basis
for the groundedness metric in \Cref{sec:groundedness-metric}, with the knowledge
graph serving as the reference. The essential methodological requirement, adopted
here, is that the verifying judge must be independent of the system that produced
the text.

\subsection{Factuality Benchmarks and Their Limits}

FactBench~\cite{bayat2025factbench} provides a dynamic benchmark for in-the-wild
factuality evaluation. The original proposal for this thesis intended to use it
both as a source for the knowledge environment and as the hallucination
evaluator. That plan was abandoned, and the reason is methodological rather than
practical: if the evaluator's facts are contained in the graph the system
retrieves from, a low hallucination rate follows by construction rather than by
merit. The evaluation in \Cref{ch:setup} therefore measures claim grounding
against the graph with a human-validated judge, keeping the evaluator's ground
truth outside the system's evidence.

\section{Comparative Summary of Agentic KGQA Systems}\label{sec:comparison-table}

\Cref{tab:mechanisms} compares the six closest systems on the mechanisms
relevant to this thesis, and \Cref{tab:reported} records their evaluation
settings and published accuracy.

\begin{table}[!tb]
  \begin{center}
    \caption{Mechanism comparison of the six closest prior systems. No system
      performs claim-level verification of the final answer against retrieved
      triples before responding.}
    \label{tab:mechanisms}
    \footnotesize
    \begin{tabular}{|l|p{38mm}|p{22mm}|p{22mm}|p{25mm}|}
      \hline
      \textbf{System} & \textbf{Exploration strategy} & \textbf{Decomposes
      query?} & \textbf{Backtracks?} & \textbf{Verifies before answering?} \\
      \hline
      ToG~\cite{sun2024think} & Training-free LLM-guided iterative beam search;
      beam width and depth both 3 & No & No --- prunes within beam, cannot
      return & No \\
      \hline
      PoG~\cite{chen2024plan} & Self-correcting adaptive planning with Guidance,
      Memory, Reflection & Yes, into sub-objectives with constraints & Yes ---
      Reflection selects the entity to return to & Partial --- corrects the
      search, not the answer \\
      \hline
      RoG~\cite{luo2024reasoning} & Planning--retrieval--reasoning over generated
      relation paths & Partial --- relation-path plans & No & Faithful by
      construction; no explicit check \\
      \hline
      KG-Agent~\cite{jiang2024kgagent} & Toolbox with KG executor and dynamic
      memory; autonomous tool selection & Implicit, via tool calls & No & No \\
      \hline
      GoG~\cite{xu2024generate} & Thinking--searching--generating; generates
      triples when the KG is incomplete & Yes & No --- expands subgraph instead
      & No \\
      \hline
      StructGPT~\cite{jiang2023structgpt} & Iterative reading-then-reasoning via
      invoke--linearise--generate interfaces & No & No & No \\
      \hline
      KBQA-o1~\cite{luo2025kbqao1} & Monte Carlo Tree Search with policy and
      reward models over a ReAct agent; fine-tuned & Implicit, stepwise logical
      form & Via MCTS tree revisiting & No \\
      \hline
      \textbf{AGR} (this work) & Guided best-first beam search with explicit
      backtracking & Yes, into ordered sub-objectives & Yes --- threshold, dead
      end, or evaluator veto & \textbf{Yes --- claim-level, pre-generation} \\
      \hline
    \end{tabular}
  \end{center}
\end{table}

\begin{table}[!tb]
  \begin{center}
    \caption{Reported evaluation settings and Hits@1 as published. These figures
      are \emph{not} directly comparable to one another: backbones differ, some
      papers evaluate on sampled subsets, and the backbone accounts for much of
      the variation.}
    \label{tab:reported}
    \footnotesize
    \begin{tabular}{|l|l|p{45mm}|c|c|}
      \hline
      \textbf{System} & \textbf{Backbone} & \textbf{Datasets} &
      \textbf{WebQSP} & \textbf{CWQ} \\
      \hline
      ToG & GPT-3.5 & CWQ, WebQSP, GrailQA~\cite{gu2021beyond}, QALD10-en,
      SimpleQuestions, WebQuestions, T-REx, Zero-Shot RE, Creak & 76.2 & 58.9 \\
      \cline{2-5}
       & GPT-4 & & 82.6 & 69.5 \\
      \hline
      PoG & GPT-3.5 & CWQ, WebQSP, GrailQA & $\approx$82.0 & $\approx$63.2 \\
      \cline{2-5}
       & GPT-4 & & $\approx$87.3 & $\approx$75.0 \\
      \hline
      RoG & Fine-tuned LLaMA2-7B & WebQSP, CWQ & 85.7 & 62.6 \\
      \hline
      KG-Agent & Fine-tuned 7B & WebQSP, CWQ, plus other KBs & 83.3 & 72.2 \\
      \hline
      GoG & GPT-3.5 & WebQSP, CWQ (complete and incomplete KG) & 78.7 & 55.7 \\
      \cline{2-5}
       & GPT-4 & & 84.4 & 75.2 \\
      \hline
      StructGPT & GPT-3.5 & WebQSP, MetaQA, TableQA, Text-to-SQL & 72.6 & --- \\
      \hline
      KBQA-o1 & Llama-3.1-8B & GrailQA, WebQSP, GraphQ; 100-shot low-resource
      protocol & 59.8$^{\dagger}$ & --- \\
      \cline{2-5}
       & Llama-3.3-70B & & 67.0$^{\dagger}$ & --- \\
      \hline
    \end{tabular}

    \vspace{1mm}
    {\footnotesize $^{\dagger}$ KBQA-o1 reports \textbf{F1}, not Hits@1, and
    evaluates under a 100-shot low-resource protocol; these cells are therefore
    \emph{not} comparable with the rest of the column and are shown only to
    situate the system. It does not evaluate CWQ.}
  \end{center}
\end{table}

Three caveats attach to \Cref{tab:reported} and are carried into every
comparison in this thesis.

First, the figures are drawn from the original publications and reflect different
backbones. RoG and KG-Agent are fine-tuned, and their numbers are not comparable
with prompting methods at all. Among the prompting methods, the backbone accounts
for much of the spread --- the same methods drop sharply when run on smaller
open-weight models. This is the direct motivation for holding the backbone
constant across every system in \Cref{ch:setup}: the literature does not, so a
cross-paper table cannot settle which architecture is better.

Second, evaluation protocols differ. Some papers evaluate on sampled test
subsets, and answer-matching conventions are not uniform across implementations.

Third, the Plan-on-Graph figures are marked as approximate. They are consistent
with values reported in follow-up work, but should be confirmed against the
primary tables before being relied upon in any quantitative claim. No claim in
this thesis depends on them; they appear here to situate the design space.

\section{Positioning of This Work}\label{sec:positioning}

\Cref{tab:mechanisms} makes the gap explicit. Decomposition is well established
--- PoG and GoG both do it. Backtracking exists, in PoG. Faithful grounding by
construction exists, in RoG. But in the final column, \emph{no surveyed system
decomposes its draft answer into atomic claims and checks each against the
retrieved triples before responding}. PoG's Reflection corrects the search while
it is running; RoG's answers are grounded because of how they are produced, not
because anything examined them. The slot is open.

AGR's contribution is therefore positioned as follows.

The \textbf{primary} contribution is the Structural Verification Layer: a
claim-level check against traversed triples, executed before the answer is
emitted, with repair policies that either resume exploration or remove
unsupported content. \Cref{ch:verification} specifies it.

The \textbf{secondary} contribution is a systematic component ablation
quantifying what each agentic mechanism actually contributes, in accuracy and in
tokens. As \Cref{tab:mechanisms} shows, these mechanisms are typically reported
as bundles; which of them earns its cost is not established in this literature.

A deliberate consequence of this positioning is that the thesis does not depend
on beating prior systems on raw Hits@1. The comparison that matters is against
the same architecture with the verification layer removed, on the same graph,
with the same backbone --- which is why the ablation in \Cref{sec:ablations},
rather than \Cref{tab:reported}, carries the argument. Where AGR is compared
against a prior system, that comparison is run under controlled conditions in
this thesis's own environment rather than quoted across papers.

\endinput
